%! AP Statistics — Unit 7, Lesson 7.1 — Homework (Answer Key)
\documentclass[10pt]{article}
\usepackage{apstats-article}
\usepackage{apstats-key}

\begin{document}

\pageheader{Unit 7, Lesson 7.1}{Homework --- Answer Key}
\namedateperiod

\vspace{0.1in}
\textbf{Directions:} For each scenario, complete all parts. Show your reasoning clearly.

\begin{enumerate}

  \item \textbf{Airport Security.}
        ($H_0{:}\ \mu = 15$ min, $H_a{:}\ \mu > 15$ min, $\alpha = 0.05$)

        \begin{enumerate}[label=\alph*.]
          \item Describe a Type I error in context.
                \ansline{Concluding that the mean wait time exceeds 15 minutes when it actually equals 15 minutes --- resources are spent on improvements that were not needed.}
          \item Describe a Type II error in context.
                \ansline{Failing to conclude that the mean wait time exceeds 15 minutes when it actually does --- the agency takes no action while travelers endure excessive wait times.}
          \item Which error is more consequential for travelers?
                \ansline{Type II: travelers continue to experience long waits without any corrective action. Type I wastes agency resources but does not directly harm travelers.}
        \end{enumerate}

  \item \textbf{Online Learning Platform.}
        ($H_0{:}\ \mu = 72$, $H_a{:}\ \mu \ne 72$, $\alpha = 0.05$)

        \begin{enumerate}[label=\alph*.]
          \item Describe a Type I error in context.
                \ansline{Concluding that the platform changes mean scores when it actually has no effect --- the district may expand or abandon a neutral program based on a false signal.}
          \item Describe a Type II error in context.
                \ansline{Failing to detect a real effect of the platform on mean scores --- a beneficial or harmful platform goes unrecognized.}
          \item Effect of lowering $\alpha$ to 0.01:
                \ansline{The Type I error rate decreases from 5\% to 1\% --- the district is less likely to make a decision based on a false positive. However, the Type II error rate increases --- real effects of the platform are more likely to go undetected. The test requires stronger evidence before concluding the platform makes a difference.}
        \end{enumerate}

  \item \textbf{Environmental Testing.}
        ($H_0{:}\ \mu = 10$ ppb, $H_a{:}\ \mu < 10$ ppb, $\alpha = 0.05$)

        \begin{enumerate}[label=\alph*.]
          \item Both error types:
                \ansline{Type I: concluding lead levels are below 10 ppb when they are actually at 10 ppb --- the facility declares water safe without sufficient evidence. Type II: failing to conclude lead levels are below 10 ppb when they actually are --- the facility continues expensive treatment unnecessarily.}
          \item More dangerous error:
                \begin{teachernote}
                    Note: the alternative here is $\mu < 10$ ppb, meaning the facility
                    \emph{wants} to show the water is safe. A Type I error (incorrectly
                    declaring water safe) is the more dangerous error for public health.
                \end{teachernote}
                \ansline{Type I: incorrectly concluding that lead levels are below the safe threshold when they are not puts the public at risk of lead exposure.}
        \end{enumerate}

\end{enumerate}

\begin{reflectionbox}
\textbf{Reflection:}
\ansline{Reducing $\alpha$ (lowering the Type I error rate) means requiring stronger evidence to reject $H_0$. This makes it harder to reject $H_0$ even when it is false, increasing the Type II error rate. The only way to decrease both error rates simultaneously is to increase the sample size, which reduces variability and makes the test more sensitive.}
\end{reflectionbox}

\end{document}
